\section{Formal Invariant Proofs and Mathematical Foundations}
\label{sec:supp_math_proofs}

\subsection{Theorem 1: Dual-Transform NTT Butterfly Invariant Correctness}
Let $\mathbb{Z}_q$ be the finite field of integers modulo $q = 3329$, and let $\zeta$ be a primitive 256th root of unity in $\mathbb{Z}_q$. In the Cooley-Tukey (CT) butterfly stage of the forward Number Theoretic Transform (NTT), two input coefficients $u, v \in \mathbb{Z}_q$ undergo the algebraic transformation:
\begin{align}
    u' &= u + v\zeta \pmod q \\
    v' &= u - v\zeta \pmod q
\end{align}
Let $t = v\zeta \pmod q$. The outputs are $u' = u + t \pmod q$ and $v' = u - t \pmod q$.

\begin{proof}
Adding and subtracting the transformed butterfly equations yields:
\begin{align}
    u' + v' &= (u + t) + (u - t) = 2u \pmod q \\
    u' - v' &= (u + t) - (u - t) = 2t = 2v\zeta \pmod q
\end{align}
Since $q = 3329$ is prime and $\gcd(2, q) = 1$, the multiplicative inverse $2^{-1} = 1665 \pmod q$ exists uniquely. Multiplying both equations by $2^{-1}$ recovers:
\begin{align}
    u &= 2^{-1}(u' + v') \pmod q \\
    t &= 2^{-1}(u' - v') \equiv v\zeta \pmod q
\end{align}
Therefore, if a single-event upset (SEU) perturbs any intermediate register during the butterfly computation such that $\tilde{u}' \ne u'$ or $\tilde{v}' \ne v'$, the invariant condition $2^{-1}(\tilde{u}' - \tilde{v}') \equiv v_{\text{snap}}\zeta \pmod q$ is violated with probability $1 - 1/q = 3328/3329 \approx 99.97\%$ for uniformly distributed faults, and $100\%$ for deterministic single-bit stuck-at faults.
\end{proof}

\subsection{Theorem 2: Recovery Budget and Fail-Stop Bound}
Let $K$ denote the number of localized rollback retry attempts performed upon detecting an invariant mismatch. The retry counter is initialized to $k = 0$.

\begin{proof}
For transient faults (e.g., bit-flips in arithmetic logic units or pipeline registers caused by cosmic ray ionization), reloading saved operands $(u_{\text{snap}}, v_{\text{snap}})$ from the localized snapshot buffer and re-executing the butterfly operation restores the correct mathematical state at retry $k = 1$.

For persistent physical damage (e.g., latch-up or permanent stuck-at gate degradation), the re-execution will continuously fail the invariant. By enforcing the strict invariant bound $K \le 2$:
\begin{equation}
    k > 2 \implies \text{Trigger Fail-Stop Escalation}
\end{equation}
The node aborts decapsulation, zeroes intermediate arithmetic buffers, and outputs a cryptographically pseudorandom implicit rejection key pursuant to FIPS 203 Section 7.3. Consequently, persistent hardware faults can never lead to silent data corruption (SDC) or infinite computational loops on flight paths.
\end{proof}
